%! Scatterplots and Correlation — Unit 2, Lesson 2.2 — Guided Notes & Practice (Answer Key)
% MIRROR of ../notes/main.tex: -key in place of -boxes, \termblank -> \termans, \blank -> \ans, and the
% fourth argument of every \pcell filled. Nothing else differs — keep it that way.
% NO teachernote here — teacher prose lives in the lesson plan.
% THE in-class document, in the MAIN IDEAS / NOTES shape (LESSON_SHAPE.md §1), at the COMPACT
% budget (2026-09-29): the vocabulary box, then ONE guidednotes table of four rows, 13 problems,
% two pages. Each row is one idea — a label on the left; on the right a terse definition the
% student completes, then a grid of short numbered problems with reserved answer space.
%   I DO   : the teacher fills the definition lines and works the FIRST problem of each grid
%            (1, 4, 7).
%   WE DO  : the class works the rest of each grid; the last row, Guided Practice (11–13), is
%            worked entirely together in a fresh context.
%   YOU DO : nothing on this page — ../homework/ IS the individual practice.
% TRAP  : problem 6 — the heaviest car is extreme in x but FOLLOWS the pattern; the hybrid is the
%         unusual point (EK DAT-1.A.6).
% CRUX  : problem 9 — r = 0.03 for speed vs. mpg, and the pattern is a strong CURVE. r near 0
%         means no LINEAR association, not no relationship (EK DAT-1.B.1, DAT-1.C.1).
%
% THE DATA
%   Cars (weight, 1000 lb / highway mpg), 12 models, r = -0.86:
%     2.5/38 2.8/36 3.0/35 3.2/48 (hybrid) 3.3/32 3.5/31 3.6/29 3.9/28 4.1/26 4.4/24 4.8/21 5.4/17
%   One car driven at 12 speeds (mph / mpg), r = 0.03:
%     20/23 25/27 30/30 35/32 40/34 45/35 50/35 55/34 60/32 65/30 70/27 75/24
%   Pool (high temp °F / visitors), 14 days, r = 0.90:
%     72/140 75/120 77/170 79/150 80/220 82/160 84/200 85/250 87/210 89/280 90/230 92/300
%     95/260 98/330
%
% NO SKETCH QUESTIONS: every scatterplot is pre-drawn in TikZ and read.
\documentclass[12pt]{article}
\usepackage{apstats-article}
\usepackage{apstats-key}   % pulls in -boxes; do NOT also load -boxes

\begin{document}

\pageheader{Unit 2, Lesson 2.2}{Guided Notes \& Practice --- Answer Key}
% NAMESTRIP: no \namedateperiod here — the name row lives on the cover only.

\vspace{2pt}

\boxguard[16]
\begin{vocabbox}
\small
Fill in each term \textbf{as we name it} in the notes below.
\par
\termans{Scatterplot}{a graph of two quantitative variables on the same individuals; one dot per individual.}
\par
\termans{Explanatory \& response}{the explanatory variable ($x$-axis) helps explain or predict the response ($y$-axis), the outcome.}
\par
\termans{Correlation $r$}{a number from $-1$ to $1$ giving the direction and strength of a LINEAR association.}
\par
\end{vocabbox}

\small
\begin{guidednotes}
% ── ROW 1: scatterplots, explanatory and response ────────────────────────────
\mainidea[Two quantitative]{Variables} &
\begin{minipage}[c]{0.49\linewidth}\raggedright
A \textbf{scatterplot} shows two \ans{quantitative} variables on the same individuals,
one dot per \ans{individual}.\par\vspace{3pt}
The \textbf{explanatory} variable (\ans{$x$}-axis) helps \ans{explain/predict} the
\textbf{response} (\ans{$y$}-axis).
\end{minipage}\hfill
\begin{minipage}[c]{0.49\linewidth}\centering
\begin{tikzpicture}[x=1.2cm, y=0.075cm]
  \foreach \y in {20,30,40,50} {
    \draw[linegray, very thin] (2,\y) -- (6,\y);
    \node[left, font=\tiny] at (2,\y) {\y};}
  \draw[thick] (2,15) -- (6.1,15);
  \draw[thick] (2,15) -- (2,52);
  \foreach \x in {2,3,4,5,6} {
    \draw (\x,14.3) -- (\x,15.7); \node[below, font=\tiny] at (\x,14.5) {\x};}
  \foreach \w/\m in {2.5/38, 2.8/36, 3.0/35, 3.2/48, 3.3/32, 3.5/31, 3.6/29, 3.9/28,
                     4.1/26, 4.4/24, 4.8/21, 5.4/17}
    \fill[navy] (\w,\m) circle (1.6pt);
  \node[font=\tiny] at (4,6.5) {Weight (1000 lb)};
  \node[rotate=90, font=\tiny] at (1.45,33) {Highway mpg};
  \node[font=\scriptsize\bfseries] at (4,56) {12 car models};
\end{tikzpicture}
\end{minipage}
\notesprompt{Read the scatterplot.}
\begin{probgrid}{|Y|Y|Y|}
\pcell{1}{Name the explanatory and the response variable.}{1.2cm}{Expl.: weight\par Resp.: mpg} &
\pcell{2}{What does one dot stand for?}{1.2cm}{One car model} &
\pcell{3}{About how many mpg does the heaviest car get?}{1.2cm}{About 17 mpg} \\ \hline
\end{probgrid}
\\ \hline
% ── ROW 2: describing an association (DOFS, in context) ──────────────────────
\mainidea[Describing the]{Pattern} &
Describe a scatterplot with \textbf{DOFS}, always in context:\par\vspace{2pt}
\textbf{D}irection --- \textbf{positive}: as $x$ increases, $y$ tends to \ans{increase};
\textbf{negative}: $y$ tends to \ans{decrease}.\par
\textbf{O}utliers \& unusual features --- points far from the \ans{pattern}, or clusters.\par
\textbf{F}orm --- \ans{linear} (a straight-line pattern) or curved.\par
\textbf{S}trength --- how closely the points follow the \ans{form}: strong, moderate, or
weak.
\\
& \notesprompt{Describe the 12 cars above.}
\begin{probgrid}{|Y|Y|Y|}
\pcell{4}{Direction --- and what it means for these cars.}{1.5cm}{Negative: heavier cars tend to get lower mpg.} &
\pcell{5}{Form and strength.}{1.5cm}{Linear; fairly strong.} &
\pcell{6}{The heaviest car sits far right. Is it the unusual point?}{1.5cm}{No, it fits the pattern. The hybrid (48 mpg) is.} \\ \hline
\end{probgrid}
\\ \hline
% ── ROW 3: correlation — THE CRUX (problem 9) ────────────────────────────────
\mainidea[Correlation]{Reading $r$} &
\begin{minipage}[c]{0.58\linewidth}\raggedright
The \textbf{correlation} $r$ measures the direction and strength of the \ans{linear}
association between two quantitative variables. It is always between \ans{$-1$} and
\ans{$1$}; its sign gives the \ans{direction}, and closer to $\pm 1$ means a
\ans{stronger} linear pattern. $r$ has no \ans{units}, and correlation does not prove
\ans{causation}.
\end{minipage}\hfill
\begin{minipage}[c]{0.40\linewidth}\centering
\begin{tikzpicture}[x=0.068cm, y=0.15cm]
  \foreach \y in {20,25,30,35} {
    \draw[linegray, very thin] (15,\y) -- (80,\y);
    \node[left, font=\tiny] at (15,\y) {\y};}
  \draw[thick] (15,20) -- (81,20);
  \draw[thick] (15,20) -- (15,37);
  \foreach \x in {20,40,60,80} {
    \draw (\x,19.7) -- (\x,20.3); \node[below, font=\tiny] at (\x,19.8) {\x};}
  \foreach \s/\m in {20/23, 25/27, 30/30, 35/32, 40/34, 45/35, 50/35, 55/34, 60/32,
                     65/30, 70/27, 75/24}
    \fill[navy] (\s,\m) circle (1.6pt);
  \node[font=\tiny] at (47.5,15.8) {Speed (mph)};
  \node[rotate=90, font=\tiny] at (6.5,28) {mpg};
  \node[font=\scriptsize\bfseries] at (47.5,39.5) {One car, 12 speeds: $r = 0.03$};
\end{tikzpicture}
\end{minipage}
\notesprompt{Interpret $r$.}
\begin{probgrid}{|Y|Y|}
\pcell{7}{For the 12 cars, $r = -0.86$. What do its sign and size tell you?}{1.1cm}{Negative; a fairly strong linear association.} &
\pcell{8}{Weigh the cars in kilograms and swap the axes. Now what is $r$?}{1.1cm}{Still $-0.86$: no units, and $x$/$y$ order does not matter.} \\ \hline
\pcell{9}{For the speed plot, $r = 0.03$. A student says, ``So speed has no relationship with fuel economy.'' Respond.}{1.5cm}{Wrong: a strong \emph{curved} relationship. $r$ measures only \emph{linear} association.} &
\pcell{10}{In 40 countries, $r = 0.81$ between cars per 1000 people and life expectancy. Do cars make people live longer?}{1.5cm}{Not shown --- correlation is not causation; national wealth drives both.} \\ \hline
\end{probgrid}
\\ \hline
% ── ROW 4: guided practice — WE DO, together, fresh context ──────────────────
\mainidea[Guided practice]{The City Pool} &
\begin{minipage}[c]{0.45\linewidth}\raggedright
A city pool recorded the day's high temperature ($^\circ$F) and the number of visitors on 14
summer days. Technology gives $r = 0.90$.
\notesprompt{We work these together.}
\end{minipage}\hfill
\begin{minipage}[c]{0.53\linewidth}\centering
\begin{tikzpicture}[x=0.17cm, y=0.0105cm]
  \foreach \y in {100,200,300} {
    \draw[linegray, very thin] (70,\y) -- (100,\y);
    \node[left, font=\tiny] at (70,\y) {\y};}
  \draw[thick] (70,100) -- (100.5,100);
  \draw[thick] (70,100) -- (70,345);
  \foreach \x in {70,80,90,100} {
    \draw (\x,96) -- (\x,104); \node[below, font=\tiny] at (\x,97) {\x};}
  \foreach \t/\v in {72/140, 75/120, 77/170, 79/150, 80/220, 82/160, 84/200, 85/250,
                     87/210, 89/280, 90/230, 92/300, 95/260, 98/330}
    \fill[navy] (\t,\v) circle (1.6pt);
  \node[font=\tiny] at (85,50) {High temperature ($^\circ$F)};
  \node[rotate=90, font=\tiny] at (64.5,222) {Visitors};
\end{tikzpicture}
\end{minipage}
\begin{probgrid}{|Y|Y|Y|}
\pcell{11}{Name the explanatory and the response variable.}{1.6cm}{Expl.: temperature\par Resp.: visitors} &
\pcell{12}{Describe the association in context (DOFS and $r$).}{1.6cm}{Strong, positive, linear: hotter days tend to bring more visitors.} &
\pcell{13}{Snack-bar ice-cream sales and visitors: $r = 0.88$. Does ice cream bring swimmers?}{1.6cm}{No --- hot weather drives both.} \\ \hline
\end{probgrid}
\\ \hline
\end{guidednotes}

% ── THE NOTES END AT GUIDED PRACTICE ─────────────────────────────────────────
% There is deliberately no "you do" here: ../homework/ is the individual practice,
% started in class in the last 8 minutes while the teacher circulates.

\end{document}
